\documentclass[a4paper,10pt]{article}
\usepackage[utf8]{inputenc}
\usepackage{enumitem}
\usepackage{geometry}
\usepackage{hyperref}
\usepackage{graphicx}

\geometry{left=2cm, right=2cm, top=2cm, bottom=2cm}

\begin{document}

\begin{center}
    {\LARGE \textbf{Igor Marcel Moreira Xavier}} \\
    \vspace{2mm}
    \href{mailto:igor2mxavier@gmail.com}{igor2mxavier@gmail.com} \quad
    \href{tel:+5551995816448}{+55 51 99581 6448} \\
    \href{https://linkedin.com/in/igormarcelmoreira}{linkedin.com/in/igormarcelmoreira}
\end{center}

\vspace{4mm}

\section*{Resumo}

Engenheiro de Software com mais de 6 anos de experiência no desenvolvimento de aplicações corporativas para ambientes mobile, web e nuvem. Especialista em React, Angular, .NET MAUI e ASP.NET Core, com experiência na liderança de decisões de arquitetura, mentoria de desenvolvedores, CI/CD e implantação de soluções mobile corporativas. Recentemente, concluiu um intercâmbio acadêmico na Hankuk University of Foreign Studies, na Coreia do Sul, aprimorando habilidades de comunicação e colaboração multicultural.

\section*{Experiência Profissional}

\subsection*{TRUE --- Tecnologia para a vida (2022 -- 2025)}
\begin{itemize}[leftmargin=1cm]
    \item \textbf{Desenvolvedor Full-Stack Sênior \& Referência Técnica Mobile (2024)}
    \item Atuação como referência técnica para equipes de desenvolvimento mobile e web, orientando arquitetura, revisando código e definindo padrões em projetos simultâneos.
    \item Desenvolvimento e manutenção de pipelines CI/CD para aplicações Android e iOS, aumentando a confiabilidade das entregas e reduzindo o esforço operacional de releases.
    \item Liderança na distribuição corporativa de aplicativos via Microsoft Intune (MDM), incluindo gerenciamento de configurações e controle de acesso alinhado a RBAC.
    \item Integração da Play Integrity API para segurança em tempo de execução e verificação de autenticidade em dispositivos Android.
    \item Liderança na migração completa de tecnologia de Xamarin para .NET MAUI em 6 aplicações em produção --- gerenciando atualização de dependências, versionamento de APIs e testes de regressão.
    \item Desenvolvimento full-stack utilizando Angular/TypeScript (frontend), WCF/VB.Net/Asp.Net Core (backend) e C\# (mobile).
    \item \textbf{Projetos em destaque:}
    \begin{itemize}
        \item \textbf{CAIRHOS:} Ferramenta web para reguladores médicos visualizarem disponibilidade de unidades de saúde; liderança no desenvolvimento e coordenação da equipe.
        \item \textbf{Suite de Aplicativos Mobile (6 apps):} CHAMAR 192, SAPH Mobile, TRUE PCR, TRUE Checklist, TRUE SAPH Gestão, Unimed POA Chamar SOS --- ciclo completo desde o desenvolvimento até o deploy CI/CD para Android e iOS.
    \end{itemize}
\end{itemize}

\subsection*{Interanet IT (2020 -- 2022)}
\begin{itemize}[leftmargin=1cm]
    \item \textbf{Desenvolvedor Full-Stack}
    \item Desenvolvimento de soluções web, backend e mobile para clientes dos setores de saúde, finanças, educação e energia.
    \item \textbf{Projetos em destaque:}
    \begin{itemize}
        \item \textbf{Banrisul:} Implantação e administração do Moodle LMS na intranet corporativa para treinamento de funcionários.
        \item \textbf{Sulgás:} Manutenção de CMS Joomla, integração com APIs ERP e desenvolvimento de aplicativo mobile em Xamarin Forms.
        \item \textbf{Marista:} Desenvolvimento e manutenção em Asp.Net Core e SharePoint; automação de fluxos de trabalho.
        \item \textbf{Sinos Tecnologia:} Desenvolvimento de sistema ERP completo utilizando Angular e Asp.Net Core C\#.
        \item \textbf{Aucon:} Desenvolvimento de aplicação mobile em Xamarin Forms.
    \end{itemize}
\end{itemize}

\section*{Formação Acadêmica}

\begin{itemize}[leftmargin=2cm]
    \item[2026] \textbf{Intercâmbio Acadêmico --- Engenharia de Software}, \\
    Hankuk University of Foreign Studies (HUFS) --- Coreia do Sul
    \item[2023--Atual] \textbf{Bacharelado em Engenharia de Software}, \\
    Pontifícia Universidade Católica do Rio Grande do Sul (PUCRS) --- Brasil
    \item[2020--2022] \textbf{Engenharia de Controle e Automação (incompleto)}, \\
    Universidade Estadual do Rio Grande do Sul (UERGS) --- Brasil
\end{itemize}

\section*{Projetos Acadêmicos}

\subsection*{Lino (Hackathon 2025) --- \textit{1º Lugar}}
\begin{itemize}[leftmargin=1cm]
    \item \textbf{Função:} Líder Front-End
    \item Desenvolvimento completo da interface em React/TypeScript sob restrições de tempo; definição de padrões de código e liderança da integração frontend-backend.
\end{itemize}

\subsection*{CP-Planta (AGES 2024/2)}
\begin{itemize}[leftmargin=1cm]
    \item \textbf{Função:} Referência em Front-End e Arquitetura de Dados
    \item Desenvolvimento de interface em React com Figma; definição de fluxos de dados, modelos lógico-conceituais e contratos de integração de APIs.
\end{itemize}

\subsection*{Polymathech (AGES 2024/1) --- \textit{Projeto Destaque}}
\begin{itemize}[leftmargin=1cm]
    \item \textbf{Função:} Referência Front-End
    \item Aplicação web em React TypeScript; padronização de código e consistência de ferramentas em toda a equipe.
\end{itemize}

\subsection*{Moving The Cities: Santiago, Chile (2025)}
\begin{itemize}[leftmargin=1cm]
    \item Projeto internacional multidisciplinar voltado à resolução de desafios reais de negócios com foco em impacto social; fortalecimento de habilidades de colaboração intercultural.
\end{itemize}

\section*{Competências Técnicas}

\begin{itemize}[leftmargin=1cm]
    \item \textbf{Linguagens:} C\#, TypeScript, JavaScript, Python, HTML, CSS, Java, SQL
    \item \textbf{Frameworks \& Bibliotecas:} React, Angular, Asp.Net Core, .NET MAUI, Xamarin, Node.js, NestJS
    \item \textbf{Mobile:} Desenvolvimento Android \& iOS (MAUI, Xamarin Forms); distribuição MDM via Intune
    \item \textbf{CI/CD \& DevOps:} GitHub Actions, Azure DevOps, GitLab CI, Bitbucket Pipelines, TFS
    \item \textbf{Cloud \& Plataformas:} AWS, Firebase, Microsoft Azure
    \item \textbf{Ferramentas:} Git, VS Code, Visual Studio, Figma, Moodle, Joomla, SharePoint
    \item \textbf{Práticas:} UI/UX Design, Padrões de Código, Layouts Responsivos, Migração de Tecnologias, Integração de Sistemas
\end{itemize}

\section*{Idiomas}
\begin{itemize}[leftmargin=1cm]
    \item \textbf{Português:} Nativo
    \item \textbf{Inglês:} Avançado
\end{itemize}

\end{document}